
\section{Discussion}

\paragraph{Agreement among methods is not evidence of faithfulness.}
SHAP, Integrated Gradients, and Gradient$\times$Input form a near-identical cluster once they are aggregated the same way (MLP arm: $\rho = 0.98$--$0.995$; Table~\ref{tab:mlp-agreement}). In the linear arm they are analytically identical, so agreement there is guaranteed by construction. This cluster is only modestly related to causal effects. Importance faithfulness is statistically reliable ($\rho_{\text{imp}} = 0.42$--$0.46$, surviving every family-wise correction we tried) but explains only about 18--21\% of the variance in feature effects, and directional faithfulness is near zero and nonsignificant for all three methods. Practitioners should therefore not treat agreement among attribution methods as evidence that they identify causally important SAE features.

\paragraph{Attribution adds little over simple probe-based scores.}
The probe-side controls (Section~\ref{sec:causal-results}) show that SHAP importance does not beat a baseline with no background: the product $w_j \bar{x}_j$ reaches $\rho_{\text{imp}} = 0.47$, against $0.46$ for mean $|\phi|$. For direction, the raw probe weight ($\rho = +0.31$) and $w_j \bar{x}_j$ ($+0.35$) both outperform the signed mean $\Phi$ ($-0.11$). In the linear arm, $\Phi$ is anti-correlated with the probe's own steering response ($\rho = -0.28$ to $-0.37$), whereas $w_j \bar{x}_j$ reaches $+0.95$. The directional failure is therefore attribution error relative to the probe, produced by centering on a background mean, and does not require any probe--LM mismatch. These results concern signed-mean aggregation of background-centered attributions. They do not show that all attribution designs fail.

\paragraph{Local versus global attribution depends on the regime.}
Local SHAP ranks features by effect better than the global score only in one setting: the MLP arm, with inactive features included (Block~A, gap $+0.389$, $q = 0.008$). That result is fragile, since rescoring the original candidate sets gives $+0.222$ ($p = 0.054$), and it largely reflects local attribution knowing which features are active, because ablating an inactive feature moves the logit by exactly zero. Among active features (Block~B), the point estimates favor the global score ($-0.305$ linear, $-0.213$ MLP), though only the linear gap is nominally significant and it does not survive correction. No directional gap between local and global attribution is significant in either block. We therefore find no evidence that local attribution is more faithful than global attribution once feature activity is controlled, and we do not recommend one over the other on the basis of these data.

\paragraph{Attribution does not recover changed features beyond feature activity.}
In minimal pairs, the language model's readout tracks the known edit direction (sign accuracy $0.76$--$0.96$), so the harness is sensitive. Attribution nonetheless does not identify the changed features. Local SHAP is modestly above chance against random distractors (AUROC $0.49$--$0.62$), but an attribution-free baseline, whether a feature is active at the edited last token, reaches $0.79$--$0.85$. Once distractors are activity-matched, local and global SHAP fall below chance ($0.37$--$0.49$) in every seed. The one consistent pattern is positional: appended insertion (AUROC $0.622$) outperforms interior insertion ($0.492$), positive in 4 of 4 seeds, which suggests attribution read at the final token is more informative about edits adjacent to it. Causal recovery among active candidates is positive for both local and global SHAP ($\rho = 0.11$--$0.38$), but local SHAP's edge is absent or reversed except for appended insertion. We read this as weak evidence that $|\text{attribution}|$ tracks ablation effect, not that attribution identifies the changed features.

\paragraph{The task is largely lexical.}
A full-vocabulary bag-of-words regression reaches $0.794$--$0.806$ held-out accuracy, within $0.03$ of the SAE probe ($0.810$; Table~\ref{tab:lexical}), and bag-of-words explains $0.39$ of the probe logit's variance. SST-2 sentiment is thus largely recoverable from word presence, which leaves interaction-aware attribution little room to differ from simpler methods. This limits what these experiments can show, and it means the near-identical behavior of SHAP, IG, and GA may be partly a property of the task.

\subsection{Limitations}
\begin{itemize}
    \item \textbf{Probe--LM confound.} Attributions explain a trained probe over SAE activations, while faithfulness is scored against the language model's logit difference. Our estimates combine attribution error with any mismatch between what the probe encodes and what the model uses. The probe-side controls separate the two only in the linear arm, and the probe's effect predicts the measured effect at only $\rho \approx 0.3$--$0.4$, so the language-model test has limited resolution.
    \item \textbf{Method equivalence in the linear arm.} SHAP, IG, and GA are analytically identical for linear probes. The MLP arm is only weakly nonlinear (linear fit to its logit $R^2 = 0.99$) and uses approximate DeepSHAP, so it adds little beyond the linear control.
    \item \textbf{Limited scope.} All results come from one model (GPT-2 Small), one SAE and layer (layer~7), and one binary sentiment task (SST-2). We make no claim that the patterns, including the local--global comparison and the position effect in minimal pairs, hold elsewhere.
    \item \textbf{Statistical power and post hoc choices.} Most per-seed $k = 20$ associations do not survive pooled correction, so inference rests on the $k = 114$ run. The mean-$|\phi|$ importance scoring was adopted after the final results were seen, and the wider test families were added post hoc (Appendix~\ref{app:mc}). Significant effects should be read as suggestive.
    \item \textbf{Untested explanation.} Our account of the global directional failure is supported by the linear-arm probe-side controls, but we did not verify it in the MLP arm, which has no closed-form probe response.
    \item \textbf{Steering as causal reference.} Additive steering and ablation in SAE space may move activations off-distribution, and the steering strength $\alpha$ was calibrated against a different readout (a layer-11 probe) than the one used for evaluation. Intervention effects are a causal reference, not ground truth. Only the minimal-pair edits have known ground truth, and they cover a narrow class of sentiment-word changes.
\end{itemize}

\subsection{Future Directions}
The most informative next step is to separate attribution error from probe--model mismatch, either by attributing directly to the language model's logit difference or by intervening on the probe's output and comparing against the language model's response. Attribution designs that avoid centering on a background mean, or that are scored against a probe's own response in the nonlinear case, would test whether the directional failure is specific to signed-mean aggregation. Because the decomposition-and-control methodology is not specific to SHAP, it could be applied to other attribution methods, models, layers, SAE architectures, and tasks that are less lexically solvable than SST-2, with larger feature samples to improve power. Finally, a broader set of edits with known ground truth would test whether the position effect in minimal pairs is a general property of attribution or specific to this setup.



\section{Conclusion}
We evaluated SHAP, Integrated Gradients, and Gradient$\times$Input on SAE latent features in a probe-mediated pipeline, comparing their attributions with the effects of feature ablation and steering on the language model’s logit difference. In our GPT-2 Small, layer-7, SST-2 setting, importance faithfulness was statistically significant but modest and no better than a simple background-free baseline, while directional faithfulness was near zero. Probe-side controls in the linear arm traced the directional failure to centering attributions on a background mean. Although the three methods agreed closely with one another, this agreement did not imply correspondence with causal effects. Local attribution outperformed global averaging on importance only in the MLP arm when inactive features were included; among active features, we found no advantage. In minimal pairs with known ground truth, attribution did not identify changed features beyond what feature activity alone provided. These findings show that, in this experimental setting, agreement between attribution methods is insufficient evidence of causal faithfulness, and that importance and directional faithfulness can yield different assessments of an attribution method. Practitioners should therefore evaluate these properties separately and validate probe-mediated attributions against interventions before using them to make claims about model behavior. Because our design combines attribution error with potential probe–model mismatch, future work should separate these sources of error—for example, by comparing attributions with interventions on the probe itself—and test the robustness of these results across models, layers, and tasks.

\section{Use of AI Tools}
The authors used Claude Sonnet 5 and Opus 5 to assist with generating portions of the experimental code and with editing the paper’s prose. The authors reviewed, tested, and verified all AI-assisted code and text, and take full responsibility for the experiments, results, and conclusions reported here.